%% Font size %%
\documentclass[11pt]{article}

%% Load the custom package
\usepackage{Mathdoc}

%% Numéro de séquence %% Titre de la séquence %%
\renewcommand{\centerhead}{}

%% Spacing commands %%
\renewcommand{\baselinestretch}{1} \setlength{\parindent}{0pt}

\begin{document}

\phantom{0}
\vspace{-.5cm}

\entetedevoirs{20}


\begin{exercicedevoir}[4]
  \begin{multicols}{2}
    \fbox{
      \begin{minipage}{0.9\linewidth}
        \textbf{Programme A}
        \begin{itemize}
        \item Choisir un nombre
        \item Multiplier par 9
        \item Ajouter 6
        \item Soustraire le nombre de départ
        \end{itemize}
      \end{minipage}
    }
\newcolumn

    \begin{enumerate}
    \item Tester ce programme avec le nombre $4$. \\ \encart{3cm} 
    \item Si on note $a$ le nombre de départ, quel est le résultat du programme de calcul ? \\ \encart{3cm}
    \end{enumerate}
  \end{multicols}
\end{exercicedevoir}

\begin{exercicedevoir}[4]
  \begin{multicols}{2}
    \fbox{
      \begin{minipage}{0.9\linewidth}
        \textbf{Programme B}
        \begin{itemize}
        \item Choisir un nombre
        \item Multiplier par 6
        \item Ajouter 10
        \item Multiplier par 5
        \item Enlever 4
        \end{itemize}
      \end{minipage}
    }
\newcolumn
    \begin{enumerate}
    \item Tester ce programme avec le nombre $2$. \\ \encart{3cm}
    \item Si on note $b$ le nombre de départ, quel est le résultat du programme de calcul ? \\ \encart{3cm}
    \end{enumerate}
  \end{multicols}
\end{exercicedevoir}

\newpage
\begin{exercicedevoir}[8]

On considère les programmes de calcul suivants :

\medskip

\begin{multicols}{2}

\fbox{
\begin{minipage}{0.9\linewidth}
\textbf{Programme A}
\begin{itemize}
    \item Choisir un nombre
    \item Le multiplier par $-6$
    \item Ajouter $8$
\end{itemize}
\end{minipage}
}

\columnbreak

\fbox{
\begin{minipage}{0.9\linewidth}
\textbf{Programme B}
\begin{itemize}
    \item Choisir un nombre
    \item Soustraire $4$
    \item Le multiplier par $-2$
    \item Soustraire $4$ fois le nombre choisi
\end{itemize}
\end{minipage}
}

\end{multicols}

\begin{enumerate}[label=\alph*)]
    \item Tester ces programmes avec les nombres $-8$ et $5$. \\ \encart{4cm}
    
    \item Si on note $x$ le nombre choisi, donner l'expression littérale obtenue pour chacun des deux programmes. \\ \encart{4cm}
\end{enumerate}

\end{exercicedevoir}
\begin{exercicedevoir}[4]
Je pense à un nombre. J'ajoute $15$, puis je divise le résultat par $2$ et j'obtiens $12$. \\
Quel est ce nombre ? \\
\textit{(Toute trace de recherche sera valorisée}.) \\ \encart{4cm}
\end{exercicedevoir}
\nonewpage

\end{document}

%%% Local Variables:
%%% mode: LaTeX
%%% TeX-master: t
%%% TeX-master: t
%%% End:

